\chapter*{Assignment Details}

The following is the original assignment statement received from Admmit through NTNU's bachelor task pool, project code OPG29, before work on the project commenced. The statement was provided in Norwegian; the translation below preserves the four success criteria the brief lists.

\begin{quote}
``We will develop a new resource planner as part of our complete suite of administrative tools. The aim of the planner is to make it easy for current and new customers to plan projects in a way that secures sound project execution, optimal use of resources, and compliance with all applicable legal requirements. The tool must be intuitive and simple, must ensure that all assignments are carried out in accordance with applicable requirements, must ensure that all resources are used optimally, and must ensure that all planned projects can be executed in line with regulation. The resource planner forms part of a larger system and must be able to communicate with existing databases such as vehicle and equipment overviews and employee rosters.''
\end{quote}

The brief fixed the topic but not the empirical foundation: how Norwegian traffic coordinators actually plan a day's assignments, and what software they use today, were left to the project to investigate. \Cref{sec:defining-task} sets out how the brief was operationalised through seven semi-structured interviews with traffic coordinators and transport managers across Norway, and \Cref{sec:dsr-dsrm} maps the result onto Peffers' six DSRM activities.
